\begin{abstract}
A preserved motif count is not a functional certificate for a graph intervention.
We formalize a task-information audit for a fixed two-step connectome propagator:
compute the source-to-readout transfer $T$, then test whether both task sources
remain visible. For an ideal noiseless subspace XOR task, both source columns
being nonzero is necessary and sufficient for perfect Bayes accuracy. The audit
changes three conclusions on measured MaleCNS wiring. A named-edge intervention
and a path-conditioned intervention both remove every common-relay triple and
preserve degree and module counts, yet only the first retains both sources.
Across two $k=8$ extractions, every nonzero dose retains ideal label information
while finite-readout accuracy on the nonzero doses still rises from about 67\% to above 94\%. The zero dose remains at chance. Degree and
protected nulls can also regain the task after a zero-dose graph has lost it.
Replaying the archived MLP on these reconstructions recovers the high-versus-chance split, whereas a linear readout stays at chance. On a visual parity XOR, the same MLP reaches 0.746, 0.663, and 0.506 as the transfer goes from invertible, to parallel, to a deleted source. With four CIFAR-10 images driving the operator together, a frozen embedding that itself scores 0.571 is recovered at 0.566 when the labeled source keeps a private direction. The degree-matched path-zero graph falls to 0.096 on that source and stays at 0.440 on a preserved source. A separate reverse-support CIFAR-10 subset control detects no connectome-specific
accuracy gain. The contribution is an audit that separates structural
preservation, source access, ideal label information, and finite trained behavior.
\end{abstract}
